\paragraph{二项--几何谱测度的 Fibonacci 塌缩：幂矩、Perron 固定点方程与 Hankel 乘积证书}\label{par:pom-replica-softcore-fibonacci-moment-collapse}
沿用 \( \varphi=\frac{1+\sqrt5}{2} \)、\(\psi=\frac{1-\sqrt5}{2}=-\varphi^{-1}\)。
在命题 \ref{prop:pom-multiplicity-composition-replica-explicit-eigs} 的固定谱簇上，对每个整数 \(q\ge 1\) 引入节点
\begin{equation}\label{eq:pom-replica-softcore-di-nodes}
\boxed{\;
d_i:=\frac12\,\varphi^{q-i}\psi^i=\frac12(-1)^i\varphi^{q-2i}\qquad(0\le i\le q).}
\end{equation}
并定义参数
\begin{equation}\label{eq:pom-replica-softcore-alpha-beta}
\boxed{\;
\alpha:=1+\frac{2}{\sqrt5}=\frac{\varphi^3}{\sqrt5},\qquad
\beta:=1-\frac{2}{\sqrt5}=-\frac{\psi^3}{\sqrt5},\qquad
w_i:=\binom{q}{i}\alpha^{q-i}\beta^i.}
\end{equation}

\begin{proposition}[\texorpdfstring{$S_q$}{Sq}-对称子空间上的张量谱基与二项权重闭式]\label{prop:pom-replica-softcore-sym-eigenbasis-binomial-weights}
\leanverified{paper\_pom\_replica\_softcore\_sym\_eigenbasis\_binomial\_weights}
取整数 \(q\ge 1\)，并令 \(E=\RR^2\)、\(K=\bigl(\begin{smallmatrix}1&1\\[2pt]1&0\end{smallmatrix}\bigr)\)、\(D^{(q)}:=K^{\otimes q}\)。
记 \(u:=(1,1)^{\mathsf T}\in E\) 且 \(\mathbf{1}:=u^{\otimes q}\in E^{\otimes q}\cong \RR^{U_q}\)，并令 \(\mathcal H_{\mathrm{sym}}\subseteq E^{\otimes q}\) 为 \(S_q\) 作用的不变子空间（\(\dim\mathcal H_{\mathrm{sym}}=q+1\)）。
则存在 \(\mathcal H_{\mathrm{sym}}\) 上的正交规范特征基 \(\{f_i\}_{i=0}^{q}\) 使得
\[
D^{(q)}f_i=(2d_i)f_i,
\qquad
d_i=\frac12(-1)^i\varphi^{q-2i},
\]
并且 \(\mathbf{1}\) 在该基下的投影满足
\[
\langle f_i,\mathbf{1}\rangle^2=w_i,
\qquad
\frac{w_i}{2^q}=\binom{q}{i}p^{\,q-i}(1-p)^i,
\qquad
p:=\frac{\alpha}{2}=\frac12+\frac{1}{\sqrt5}.
\]
特别地，\(\sum_{i=0}^q w_i=2^q\)，并且
\[
\prod_{i=0}^{q}w_i
=
\left(\prod_{i=0}^{q}\binom{q}{i}\right)\cdot 5^{-\frac{q(q+1)}{2}}.
\]
\end{proposition}

\begin{proof}
矩阵 \(K\) 的两特征值为 \(\varphi\) 与 \(\psi=-\varphi^{-1}\)，并可取一组正交规范特征向量 \(v_+,v_-\in E\) 使
\(
Kv_+=\varphi v_+
\)、\(
Kv_-=\psi v_-
\)、\(
v_+\perp v_-
\)。
在 \(E^{\otimes q}\) 中，张量积基向量 \(\bigotimes_{j=1}^q v_{\varepsilon_j}\)（\(\varepsilon_j\in\{+,-\}\)）为 \(D^{(q)}=K^{\otimes q}\) 的本征向量，其本征值为 \(\varphi^{q-i}\psi^i\)，其中 \(i\) 为符号 “\(-\)” 的出现次数。
将这些向量对 \(S_q\) 的置换作用作对称化并按 \(i\) 分层，即可在 \(\mathcal H_{\mathrm{sym}}\) 上得到正交规范特征基 \(\{f_i\}_{i=0}^{q}\)，其本征值为 \(\varphi^{q-i}\psi^i=(-1)^i\varphi^{q-2i}=2d_i\)。
\par
另一方面，\(\mathbf{1}=u^{\otimes q}\)。
将 \(u\) 在正交基 \(\{v_+,v_-\}\) 下展开为 \(u=\alpha_0 v_+ +\beta_0 v_-\)，则
\(
u^{\otimes q}
=\sum_{i=0}^q \sqrt{\binom{q}{i}}\alpha_0^{\,q-i}\beta_0^{\,i} f_i
\)
且
\(
w_i=\langle f_i,\mathbf{1}\rangle^2=\binom{q}{i}\alpha_0^{\,2(q-i)}\beta_0^{\,2i}
\)。
在当前归一化下有
\(
\alpha_0^{2}=\alpha
\)、\(
\beta_0^{2}=\beta
\)（见 \eqref{eq:pom-replica-softcore-alpha-beta}），从而 \(w_i=\binom{q}{i}\alpha^{q-i}\beta^i\)，与 \eqref{eq:pom-replica-softcore-alpha-beta} 一致。
由 \(\alpha+\beta=2\) 得 \(\sum_i w_i=(\alpha+\beta)^q=2^q\)。
又由 \(\alpha\beta=1/5\) 得
\[
\prod_{i=0}^{q}w_i
=\left(\prod_{i=0}^{q}\binom{q}{i}\right)\alpha^{\sum_{i}(q-i)}\beta^{\sum_i i}
=\left(\prod_{i=0}^{q}\binom{q}{i}\right)(\alpha\beta)^{\frac{q(q+1)}{2}}
=\left(\prod_{i=0}^{q}\binom{q}{i}\right)5^{-\frac{q(q+1)}{2}}.
\]
最后，令 \(p=\alpha/2\) 即得 \(w_i/2^q=\binom{q}{i}p^{q-i}(1-p)^i\)。证毕。
\end{proof}

\begin{proposition}[秩一扰动的割线方程接口（例外谱的一维刻画）]\label{prop:pom-replica-softcore-exceptional-secant-equation}
\leanverified{paper\_pom\_replica\_softcore\_exceptional\_secant\_equation}
取整数 \(q\ge 2\)，并令 \(\rho_q=\rho(A^{(q)})\) 为例外块的 Perron 根（见定理 \ref{thm:pom-replica-softcore-exceptional-spectrum-trace} 的记号）。
则 \(\rho_q\) 满足严格割线方程
\begin{equation}\label{eq:pom-replica-softcore-secant-equation}
\boxed{\;
\sum_{i=0}^{q}\frac{w_i}{\rho_q-d_i}=2.}
\end{equation}
\end{proposition}
\begin{proof}[证明要点]
由命题 \ref{prop:pom-multiplicity-composition-replica-explicit-eigs}，固定谱簇在 \((S_q)\)-对称子空间上以 \(\{d_i\}_{i=0}^{q}\) 为节点呈对角化；
而 \(T^{(q)}=\frac12(\mathbf{1}\mathbf{1}^{\mathsf T}+D^{(q)})\) 的 \(\mathbf{1}\mathbf{1}^{\mathsf T}\) 项在该子空间上给出秩一耦合。
在该谱基下应用秩一行列式引理（等价于 rank-one resolvent 恒等式），得到例外特征值 \(\lambda\) 满足
\(
1-\frac12\sum_i \frac{w_i}{\lambda-d_i}=0
\)，即 \eqref{eq:pom-replica-softcore-secant-equation}。
权重 \(w_i\) 的闭式由对称张量本征基与 \(\mathbf{1}\) 的投影系数显式计算得到，见 \eqref{eq:pom-replica-softcore-alpha-beta}。证毕。
\end{proof}

\begin{corollary}[割线方程的期望形式]\label{cor:pom-replica-softcore-secant-equation-expectation-form}
\leanverified{paper\_pom\_replica\_softcore\_secant\_equation\_expectation\_form}
取 \(q\ge 2\)，令概率权重 \(\pi_i:=w_i/2^q\) 并令随机变量 \(I\sim \mathrm{Bin}(q,1-p)\)（其中 \(p=\frac12+\frac{1}{\sqrt5}\)）。
定义几何节点随机变量 \(X:=d_I\)（节点 \(d_i\) 见 \eqref{eq:pom-replica-softcore-di-nodes}）。
则例外块 Perron 根 \(\rho_q\) 的割线方程 \eqref{eq:pom-replica-softcore-secant-equation} 等价于标量期望方程
\[
\boxed{\;
\mathbb E\!\left[\frac{1}{\rho_q-X}\right]=2^{1-q}.}
\]
\end{corollary}

\begin{proof}
由命题 \ref{prop:pom-replica-softcore-exceptional-secant-equation}，
\(
\sum_i \frac{w_i}{\rho_q-d_i}=2
\)。
两边同除以 \(2^q\) 并用 \(\pi_i=w_i/2^q\) 得
\(
\sum_i \frac{\pi_i}{\rho_q-d_i}=2^{1-q}
\)。
由 \(\pi_i=\binom{q}{i}p^{q-i}(1-p)^i\)（命题 \ref{prop:pom-replica-softcore-sym-eigenbasis-binomial-weights}）可识别 \(\pi_i\) 为 \(I\sim\mathrm{Bin}(q,1-p)\) 的质量函数，故左端为 \(\mathbb E[(\rho_q-d_I)^{-1}]\)，即所示期望形式。证毕。
\end{proof}

\begin{theorem}[二项--几何谱测度的 Fibonacci 幂矩塌缩]\label{thm:pom-replica-softcore-fibonacci-power-moment-collapse}
\leanverified{paper\_pom\_replica\_softcore\_fibonacci\_power\_moment\_collapse}
在 \eqref{eq:pom-replica-softcore-di-nodes}--\eqref{eq:pom-replica-softcore-alpha-beta} 的记号下，对任意整数 \(m\ge 0\) 成立严格恒等式
\begin{equation}\label{eq:pom-replica-softcore-moment-collapse}
\boxed{\;
\sum_{i=0}^{q} w_i\,d_i^{\,m}
\,=\,2^{-m}\,F_{m+3}^{\,q}.}
\end{equation}
因此该谱测度的全部幂矩都落在 \(\QQ\) 中，右端不含 \(\sqrt5\) 的代数延拓项。
\end{theorem}
\begin{proof}
由 \(\psi=-\varphi^{-1}\) 得 \(d_i^{\,m}=2^{-m}\varphi^{m(q-i)}\psi^{mi}\)。
于是
\[
\sum_{i=0}^{q} w_i\,d_i^{\,m}
=2^{-m}\sum_{i=0}^{q}\binom{q}{i}(\alpha\varphi^m)^{q-i}(\beta\psi^m)^i
=2^{-m}(\alpha\varphi^m+\beta\psi^m)^{q}.
\]
由 \eqref{eq:pom-replica-softcore-alpha-beta} 的 \(\alpha=\varphi^3/\sqrt5\)、\(\beta=-\psi^3/\sqrt5\)，有
\[
\alpha\varphi^m+\beta\psi^m=\frac{\varphi^{m+3}-\psi^{m+3}}{\sqrt5}=F_{m+3},
\]
代回即得 \eqref{eq:pom-replica-softcore-moment-collapse}。证毕。
\end{proof}

\begin{theorem}[Perron 特征值的纯 Fibonacci 固定点方程]\label{thm:pom-replica-softcore-perron-fibonacci-fixed-point}
\leanverified{paper\_pom\_replica\_softcore\_perron\_fibonacci\_fixed\_point}
取整数 \(q\ge 2\)，令 \(\rho_q=\rho(A^{(q)})\) 为例外块 Perron 根。
则存在唯一 \(x_q>1\) 使得
\begin{equation}\label{eq:pom-replica-softcore-rho-xq}
\boxed{\;\rho_q=2^{q-1}x_q\;}
\end{equation}
并且 \(x_q\) 满足严格收敛的固定点方程
\begin{equation}\label{eq:pom-replica-softcore-xq-fixed-point}
\boxed{\;
x_q
=1+\sum_{m=1}^{\infty}\left(\frac{F_{m+3}}{2^{m+1}}\right)^{q}x_q^{-m}.}
\end{equation}
\end{theorem}
\begin{proof}
由命题 \ref{prop:pom-multiplicity-composition-sq-symmetry-reduction}，有 \(\rho_q=\rho(T^{(q)})\)。
取 \(\mathbf{1}\in\RR^{U_q}\) 为全 \(1\) 向量作 Rayleigh 商，得
\[
\rho_q
\ge \frac{\mathbf 1^{\mathsf T}T^{(q)}\mathbf 1}{\mathbf 1^{\mathsf T}\mathbf 1}
=\frac{4^q+3^q}{2^{q+1}}
=2^{q-1}+\frac{3^q}{2^{q+1}}
>\frac{\varphi^q}{2}
=\max_{0\le i\le q} d_i.
\]
因此对每个 \(i\) 可作几何级数展开
\[
\frac{1}{\rho_q-d_i}=\frac{1}{\rho_q}\sum_{m=0}^{\infty}\left(\frac{d_i}{\rho_q}\right)^{m},
\qquad \text{且一致绝对收敛。}
\]
将之代入 \eqref{eq:pom-replica-softcore-secant-equation} 并交换求和次序得
\[
2=\frac{1}{\rho_q}\sum_{m=0}^{\infty}\rho_q^{-m}\sum_{i=0}^{q}w_i d_i^{\,m}.
\]
由定理 \ref{thm:pom-replica-softcore-fibonacci-power-moment-collapse}，
\(
\sum_i w_i d_i^{\,m}=2^{-m}F_{m+3}^q
\)，
从而
\[
2=\frac{1}{\rho_q}\sum_{m=0}^{\infty}\left(\frac{F_{m+3}}{2}\right)^{q}\left(\frac{1}{2\rho_q}\right)^m.
\]
记 \(\rho_q=2^{q-1}x_q\)，整理并移去 \(m=0\) 项（注意 \(F_3=2\)）即得 \eqref{eq:pom-replica-softcore-xq-fixed-point}。
右端关于 \(x\in(1,\infty)\) 严格递减而左端为严格递增函数 \(x\mapsto x\)，故正解唯一。证毕。
\end{proof}

\begin{corollary}[Perron 根的两项主导展开与显式余项界]\label{cor:pom-replica-softcore-rho-two-term-remainder}
\leanverified{paper\_pom\_replica\_softcore\_rho\_two\_term\_remainder}
定义
\[
b_m:=\frac{F_{m+3}}{2^{m+1}}\quad(m\ge 1),
\qquad
\varepsilon_q:=\sum_{m=1}^{\infty}b_m^{\,q}.
\]
则 \(x_q\in(1,1+\varepsilon_q)\)，并且存在余项 \(R_q\) 使
\begin{equation}\label{eq:pom-replica-softcore-rho-two-term}
\rho_q
=2^{q-1}+\frac{3^q}{2^{q+1}}+R_q,
\end{equation}
其中余项满足显式上界
\begin{equation}\label{eq:pom-replica-softcore-rho-remainder-bound}
\boxed{\;
0<R_q
<2^{q-1}\sum_{m=2}^{\infty}b_m^{\,q}
\le \frac{5^q}{2^{2q+1}}\cdot\frac{1}{1-(\varphi/2)^q}.}
\end{equation}
因此当 \(q\to\infty\) 时有严格渐近
\[
\rho_q
=2^{q-1}+\frac{3^q}{2^{q+1}}+O\!\left(\frac{5^q}{2^{2q}}\right).
\]
\end{corollary}
\begin{proof}
由定理 \ref{thm:pom-replica-softcore-perron-fibonacci-fixed-point} 的固定点方程，
\(
x_q-1=\sum_{m\ge 1}b_m^{\,q}x_q^{-m}
\)。
因 \(x_q>1\)，得 \(0<x_q-1<\sum_{m\ge 1}b_m^{\,q}=\varepsilon_q\)，从而 \(x_q\in(1,1+\varepsilon_q)\)。
再将首项 \(m=1\) 分离并乘以 \(2^{q-1}\)，用 \(b_1=F_4/2^2=3/4\) 得到 \eqref{eq:pom-replica-softcore-rho-two-term}，并且 \(R_q>0\)。
对余项上界，用 \(x_q^{-m}\le 1\) 得
\(
R_q<2^{q-1}\sum_{m\ge 2}b_m^{\,q}
\)。
又由 Fibonacci 比值单调性 \(F_{n+1}/F_n<\varphi\)（\(n\ge 1\)）推出
\(
b_{m+1}/b_m<\varphi/2
\)，
故 \(\sum_{m\ge 2}b_m^{\,q}\) 被等比级数控制并得到 \eqref{eq:pom-replica-softcore-rho-remainder-bound}。证毕。
\end{proof}

\begin{theorem}[Perron 根的 Schur--Feshbach 方程与组合--矩级数]\label{thm:pom-replica-softcore-perron-feshbach-composition-moment-series}
\leanverified{paper\_pom\_replica\_softcore\_perron\_feshbach\_composition\_moment\_series}
令 \(n:=2^q\)，并在 \(\RR^{U_q}\) 上记 \(u:=n^{-1/2}\mathbf{1}\)、\(P:=uu^{\mathsf T}\)、\(Q:=I-P\)。
令 \(D^{(q)}=K^{\otimes q}\) 且 \(E:=\frac12D^{(q)}\)，则
\[
T^{(q)}=\frac12\bigl(\mathbf{1}\mathbf{1}^{\mathsf T}+D^{(q)}\bigr)=2^{q-1}P+E.
\]
记 \(\lambda_0^{(q)}:=\rho(T^{(q)})=\rho_q\)。
定义幂矩
\[
m_k(q):=\langle u,E^k u\rangle=\frac{1}{2^{q+k}}\,F_{k+3}^{\,q}\qquad(k\ge 1),
\]
并对 \(m\ge 1\) 定义组合系数
\[
\kappa_{m}(q):=\sum_{r=1}^{m}(-1)^{r-1}\!\!\sum_{\substack{\ell_1+\cdots+\ell_r=m\\ \ell_i\ge 1}}\ \prod_{j=1}^{r} m_{\ell_j}(q),
\]
其中内层求和遍历将 \(m\) 写成正整数和的有序组合（composition）。
则 \(\lambda_0^{(q)}\) 满足严格标量方程
\[
\boxed{\;
\lambda_0^{(q)}
=2^{q-1}+m_1(q)+\sum_{s=0}^{\infty}\frac{\kappa_{s+2}(q)}{\bigl(\lambda_0^{(q)}\bigr)^{s+1}}\;}
\]
且该级数绝对收敛。
其前三个非平凡系数可写为
\[
\kappa_2(q)=m_2(q)-m_1(q)^2,\qquad
\kappa_3(q)=m_3(q)-2m_1(q)m_2(q)+m_1(q)^3,
\]
\[
\kappa_4(q)=m_4(q)-m_2(q)^2-2m_1(q)m_3(q)+3m_1(q)^2m_2(q)-m_1(q)^4,
\]
其中
\(
m_1(q)=\frac{3^q}{2^{q+1}}
\)、\(
m_2(q)=\frac{5^q}{2^{q+2}}
\)、\(
m_3(q)=\frac{8^q}{2^{q+3}}
\)、\(
m_4(q)=\frac{13^q}{2^{q+4}}
\)。
\end{theorem}

\begin{proof}
将 \(\RR^{U_q}=\mathrm{span}\{u\}\oplus u^\perp\) 分解，并在该分解下写
\[
T^{(q)}=
\begin{pmatrix}
2^{q-1}+m_1(q) & b^{\mathsf T}\\
b & C
\end{pmatrix},
\qquad
b:=QEu,\qquad C:=QEQ.
\]
对特征值 \(\lambda_0^{(q)}\) 取 Schur 补得到
\[
\lambda_0^{(q)}=2^{q-1}+m_1(q)+b^{\mathsf T}\bigl(\lambda_0^{(q)}I-C\bigr)^{-1}b.
\]
由 \(\|C\|\le \|E\|=\frac12\|D^{(q)}\|=\frac12\varphi^q\) 且 \(\lambda_0^{(q)}\ge 2^{q-1}\)，得
\(
\|C\|/\lambda_0^{(q)}\le (\varphi/2)^q<1
\)，因此
\[
(\lambda_0^{(q)}I-C)^{-1}
=\sum_{s=0}^{\infty}\frac{C^s}{(\lambda_0^{(q)})^{s+1}}
\qquad\text{（算子范数绝对收敛）}.
\]
代入得到
\[
\lambda_0^{(q)}=2^{q-1}+m_1(q)+\sum_{s=0}^{\infty}\frac{b^{\mathsf T}C^s b}{(\lambda_0^{(q)})^{s+1}}.
\]
最后，利用 \(Q=I-P\) 的展开与恒等式
\(
\langle u,XPYu\rangle=\langle u,Xu\rangle\langle u,Yu\rangle
\)，可将 \(b^{\mathsf T}C^s b=\langle u,EQ(QEQ)^sQE u\rangle\) 展开为“把 \(s+2\) 个 \(E\) 依次分割成若干块”的有序组合求和，
其符号为 \((-1)^{r-1}\)，并贡献块积 \(\prod_j \langle u,E^{\ell_j}u\rangle=\prod_j m_{\ell_j}(q)\)。
因此 \(b^{\mathsf T}C^s b=\kappa_{s+2}(q)\)，从而得到所示级数方程。
\par
幂矩闭式由 \(E^k=2^{-k}(K^k)^{\otimes q}\) 与 \(u=n^{-1/2}\mathbf{1}\) 给出：
\(
m_k(q)=2^{-q-k}\bigl(\mathbf{1}^{\mathsf T}K^k\mathbf{1}\bigr)^q
\)，而 \(\mathbf{1}^{\mathsf T}K^k\mathbf{1}=F_{k+3}\)（由 \(K^k\) 的 Fibonacci 形式）即得 \(m_k(q)=2^{-(q+k)}F_{k+3}^q\)。
系数 \(\kappa_2,\kappa_3,\kappa_4\) 由组合定义直接展开即得。证毕。
\end{proof}

\begin{proposition}[Perron 特征向量的均匀化定量界]\label{prop:pom-replica-softcore-perron-eigenvector-uniformization-angle-bound}
设 \(v^{(q)}\) 为 \(T^{(q)}\) 的 Perron 特征向量，取 \(\|v^{(q)}\|_2=1\) 且分量全正，并仍记 \(u=2^{-q/2}\mathbf{1}\)。
则有定量界
\[
\sin\angle\bigl(v^{(q)},u\bigr)
\le
\frac{\|QEu\|_2}{2^{q-1}-\|E\|}
=
\frac{\sqrt{\frac{5^q}{2^{q+2}}-\frac{9^q}{2^{2q+2}}}}{2^{q-1}-\frac12\varphi^q},
\]
从而
\[
\|v^{(q)}-u\|_2\le \sqrt2\,\sin\angle\bigl(v^{(q)},u\bigr)
\le
\frac{\sqrt{\frac{5^q}{2^{q+1}}-\frac{9^q}{2^{2q+1}}}}{2^{q-1}-\frac12\varphi^q}.
\]
\leanverified{paper\_pom\_replica\_softcore\_perron\_eigenvector\_uniformization\_angle\_bound}
\end{proposition}

\begin{proof}
由 \(T^{(q)}=2^{q-1}P+E\) 与 \(T^{(q)}v^{(q)}=\lambda_0^{(q)}v^{(q)}\)，对等式取 \(Q\) 投影得
\[
(\lambda_0^{(q)}I-QEQ)\,Qv^{(q)}=QEu\cdot \langle u,v^{(q)}\rangle.
\]
因此
\[
\|Qv^{(q)}\|_2\le \|(\lambda_0^{(q)}I-QEQ)^{-1}\|\cdot \|QEu\|_2
\le \frac{\|QEu\|_2}{\lambda_0^{(q)}-\|E\|}
\le \frac{\|QEu\|_2}{2^{q-1}-\|E\|},
\]
其中使用了 \(\|QEQ\|\le \|E\|\) 与 \(\lambda_0^{(q)}\ge 2^{q-1}\)。
又 \(\|Qv^{(q)}\|_2=\sin\angle(v^{(q)},u)\)，得第一式。
\par
对 \(\|QEu\|_2\)，注意
\(
\|Eu\|_2^2=\langle u,E^2u\rangle=m_2(q)=\frac{5^q}{2^{q+2}}
\) 且 \(\langle u,Eu\rangle=m_1(q)=\frac{3^q}{2^{q+1}}\)，从而
\[
\|QEu\|_2^2=\|Eu\|_2^2-\langle u,Eu\rangle^2
=\frac{5^q}{2^{q+2}}-\frac{9^q}{2^{2q+2}}.
\]
最后用 \(\|E\|=\frac12\|D^{(q)}\|=\frac12\varphi^q\) 即得所示显式界；并由 \(\|v-u\|_2\le \sqrt2\,\sin\angle(v,u)\) 得第二式。证毕。
\end{proof}

\begin{theorem}[Fibonacci 幂 Hankel 行列式的完全乘积公式]\label{thm:pom-replica-softcore-fibonacci-power-hankel-determinant}
\leanverified{paper\_pom\_replica\_softcore\_fibonacci\_power\_hankel\_determinant}
对任意整数 \(q\ge 1\)，定义 \((q+1)\times(q+1)\) Hankel 矩阵
\[
H_q:=\bigl(F_{i+j+3}^{\,q}\bigr)_{0\le i,j\le q}.
\]
则成立精确闭式
\begin{equation}\label{eq:pom-replica-softcore-fibpower-hankel-det}
\boxed{\;
\det(H_q)
=\left(\prod_{i=0}^{q}\binom{q}{i}\right)
\left(\prod_{k=1}^{q}F_k^{\,2(q+1-k)}\right).}
\end{equation}
右端为整数的纯乘积表达，不含任何代数数域延拓。
\end{theorem}
\begin{proof}
由 Binet 公式与 \eqref{eq:pom-replica-softcore-alpha-beta} 的 \(\alpha,\beta\) 可写
\(
F_{n+3}=\alpha\varphi^n+\beta\psi^n
\)。
从而对每个整数 \(n\ge 0\) 有指数和表示
\[
F_{n+3}^{\,q}
=\sum_{i=0}^{q}\binom{q}{i}\alpha^{q-i}\beta^i(\varphi^{q-i}\psi^i)^{n}
=\sum_{i=0}^{q}c_i r_i^{\,n},
\]
其中 \(r_i:=\varphi^{q-i}\psi^i\)、\(c_i:=\binom{q}{i}\alpha^{q-i}\beta^i\)。
令 \(V=(r_i^{\,j})_{0\le i,j\le q}\) 且 \(C=\mathrm{diag}(c_0,\dots,c_q)\)，则 \(H_q=V^{\mathsf T}CV\)，因而
\[
\det(H_q)=(\det V)^2\prod_{i=0}^{q}c_i
=\left(\prod_{0\le i<j\le q}(r_j-r_i)^2\right)\left(\prod_{i=0}^{q}c_i\right).
\]
计算 \(\prod_i c_i=(\prod_i\binom{q}{i})(\alpha\beta)^{q(q+1)/2}\) 且 \(\alpha\beta=1/5\)，得
\(
\prod_i c_i=(\prod_i\binom{q}{i})\,5^{-q(q+1)/2}
\)。
另一方面，把 \(r_i=\varphi^q(\psi/\varphi)^i\) 写为几何节点并用 Vandermonde 分解，
再用恒等式 \(1-(-\varphi^{-2})^k=\sqrt5\,F_k/\varphi^{k}\) 整理，可得
\(
\prod_{i<j}(r_j-r_i)^2
=5^{q(q+1)/2}\prod_{k=1}^{q}F_k^{2(q+1-k)}
\)；
两者的 \(5\)-因子精确抵消，得到 \eqref{eq:pom-replica-softcore-fibpower-hankel-det}。证毕。
\end{proof}

\begin{corollary}[交错几何点的 Vandermonde 平方塌缩为纯 Fibonacci 乘积]\label{cor:pom-replica-softcore-alternating-vandermonde-fibonacci-product}
\leanverified{paper\_pom\_replica\_softcore\_alternating\_vandermonde\_fibonacci\_product}
令
\(
\xi_i:=(-1)^i\varphi^{q-2i}\ (0\le i\le q)
\)。
则成立严格恒等式
\begin{equation}\label{eq:pom-replica-softcore-alternating-vandermonde}
\boxed{\;
\prod_{0\le i<j\le q}(\xi_j-\xi_i)^2
=5^{q(q+1)/2}\prod_{k=1}^{q}F_k^{2(q+1-k)}\in\ZZ.}
\end{equation}
\end{corollary}
\begin{proof}
注意 \(r_i=\varphi^{q-i}\psi^i=\xi_i\)（由 \(\psi^i=(-1)^i\varphi^{-i}\)），
故定理 \ref{thm:pom-replica-softcore-fibonacci-power-hankel-determinant} 证明中的 Vandermonde 平方闭式直接给出 \eqref{eq:pom-replica-softcore-alternating-vandermonde}。证毕。
\end{proof}

\begin{proposition}[节点多项式判别式的域范数：Fibonacci 指数分解]\label{prop:pom-replica-softcore-node-discriminant-field-norm}
令 \(K=\QQ(\sqrt5)\)，并定义节点多项式
\[
G_q(x):=\prod_{i=0}^{q}\bigl(x-\xi_i\bigr)\in K[x],
\qquad
\xi_i:=(-1)^i\varphi^{q-2i}.
\]
则其判别式 \(\mathrm{Disc}(G_q)\in K\) 的域范数满足
\[
\boxed{\;
\mathrm N_{K/\QQ}\!\bigl(\mathrm{Disc}(G_q)\bigr)
=5^{q(q+1)}\prod_{d=1}^{q}F_d^{\,4(q+1-d)}\; }.
\]
\end{proposition}

\begin{proof}
由判别式定义，
\(
\mathrm{Disc}(G_q)=\prod_{0\le i<j\le q}(\xi_j-\xi_i)^2
\)。
推论 \ref{cor:pom-replica-softcore-alternating-vandermonde-fibonacci-product} 给出该乘积在 \(K\) 中实际上属于 \(\ZZ\) 且等于
\(
5^{q(q+1)/2}\prod_{d=1}^{q}F_d^{2(q+1-d)}
\)。
因此 \(\mathrm{Disc}(G_q)\in\ZZ\subset K\)，从而
\(
\mathrm N_{K/\QQ}(\mathrm{Disc}(G_q))=\mathrm{Disc}(G_q)^2
\)，即得到所示公式。证毕。
\end{proof}

\begin{proposition}[例外特征多项式与节点多项式的结果式：域范数的完全因子化]\label{prop:pom-replica-softcore-resultant-factorization-field-norm}
令 \(K=\QQ(\sqrt5)\)，并保持节点多项式 \(G_q\) 的定义。
记 \(D^{(q)}=K^{\otimes q}\) 在 \(\mathcal H_{\mathrm{sym}}\) 上的谱为 \(\{2d_i\}_{i=0}^q\)，并令
\[
P_q(x):=\det\!\bigl(xI-(J^{(q)}+D^{(q)})|_{\mathcal H_{\mathrm{sym}}}\bigr)\in K[x],
\qquad
J^{(q)}:=\mathbf{1}\mathbf{1}^{\mathsf T}.
\]
则成立域范数完全因子化：
\[
\boxed{\;
\mathrm N_{K/\QQ}\!\bigl(\mathrm{Res}(P_q,G_q)\bigr)
=\left(\prod_{i=0}^{q}\binom{q}{i}^2\right)\cdot
\left(\prod_{d=1}^{q}F_d^{\,4(q+1-d)}\right)\; }.
\]
\end{proposition}

\begin{proof}
在命题 \ref{prop:pom-replica-softcore-sym-eigenbasis-binomial-weights} 的正交基 \(\{f_i\}\) 下，
\(
D^{(q)}|_{\mathcal H_{\mathrm{sym}}}=\mathrm{diag}(2d_0,\dots,2d_q)
\) 且 \(J^{(q)}|_{\mathcal H_{\mathrm{sym}}}=gg^{\mathsf T}\) 其中 \(g_i=\langle f_i,\mathbf{1}\rangle\)、\(g_i^2=w_i\)。
由秩一行列式引理，
\[
P_q(x)=\prod_{i=0}^{q}(x-2d_i)\left(1-\sum_{i=0}^{q}\frac{w_i}{x-2d_i}\right).
\]
因此对每个 \(i\)，在 \(x=2d_i\) 处取极限得到
\[
P_q(2d_i)=-w_i\prod_{j\ne i}(2d_i-2d_j).
\]
从而
\[
\mathrm{Res}(P_q,G_q)
=\prod_{i=0}^{q}P_q(\xi_i)
=\pm\left(\prod_{i=0}^{q}w_i\right)\cdot \mathrm{Disc}(G_q),
\]
其中 \(\xi_i=2d_i\)，而符号对域范数无影响。
对域范数，用 \(\mathrm N_{K/\QQ}(w_i)=\binom{q}{i}^2\cdot 5^{-q}\)（由命题 \ref{prop:pom-replica-softcore-sym-eigenbasis-binomial-weights} 的 \(\alpha\beta=1/5\) 与共轭交换）得
\[
\mathrm N_{K/\QQ}\!\left(\prod_{i=0}^q w_i\right)
=\frac{\prod_{i=0}^q\binom{q}{i}^2}{5^{q(q+1)}}.
\]
再由命题 \ref{prop:pom-replica-softcore-node-discriminant-field-norm} 代入
\(
\mathrm N_{K/\QQ}(\mathrm{Disc}(G_q))=5^{q(q+1)}\prod_{d=1}^q F_d^{4(q+1-d)}
\)，相乘即得结论。证毕。
\end{proof}

\begin{remark}[审计脚本]\label{rem:pom-replica-softcore-fibonacci-moment-collapse-audit-script}
式 \eqref{eq:pom-replica-softcore-moment-collapse}、\eqref{eq:pom-replica-softcore-xq-fixed-point}、\eqref{eq:pom-replica-softcore-fibpower-hankel-det} 与 \eqref{eq:pom-replica-softcore-alternating-vandermonde} 的一致性可由脚本
\texttt{scripts/exp\_pom\_replica\_softcore\_fibonacci\_moment\_collapse\_audit.py}
在小 \(q\) 上进行严格核验。
\end{remark}

\endinput
